\newpage

\section*{\hspace{-4mm} \centering Appendix}
\vspace{3mm}

\section{Supplementary Roadmap}

The appendix supports the benchmark claims in the main paper rather than serving as a storage room for extra leaderboards. It is organized into four support blocks:
\begin{itemize}[leftmargin=5mm,itemsep=1mm,topsep=1mm]
    \item \textbf{Dataset and artifact audit:} construction provenance, artifact boundary, contamination scope, P3 intervention sampling, and release contents.
    \item \textbf{Metric and statistical definitions:} uncertainty intervals, calibration references, and formal definitions for every headline and diagnostic metric.
    \item \textbf{Full results and split checks:} complete P1--P4 tables, split-resolved tables, reasoning/thinking coverage, response-control profiles, and stage-gap diagnostics.
    \item \textbf{Mechanism and cross-domain audits:} extended related work/discussion, prompts, failure cases, selection--generation isomorphism, Tier~3 controlled-prompt results, CMA-ES, and the circuit audit.
\end{itemize}

\section{Construction Audit and Artifact Availability}

The current VEH suite is not a loose aggregation of prompts. Its optimizer-facing stages are traceable to a concrete literature-extraction and admission pipeline. The upstream extraction audit covers 209 papers and 2090 field rows, with 151 papers carrying at least one issue and 77 carrying a high-risk issue. From this audit, 71 cantilever candidates are forwarded to oracle review, and the current P2 backbone admits 52 cleaned structure-review anchors. The current benchmark then packages these anchors into split-specific task banks whose manifests store source tables, input manifest hashes, seeds, and artifact SHA256 values.

\Needspace{16\baselineskip}
\begin{center}
\captionof{table}{Construction audit for the current benchmark release. The table separates source scale, split policy, and release evidence so that the benchmark can be audited as a dataset rather than only as a set of scores.}
\label{tab:appendix_construction_audit}
\centering
\scriptsize
\setlength{\tabcolsep}{3pt}
\renewcommand{\arraystretch}{1.08}
\begin{tabularx}{\textwidth}{p{17mm}p{24mm}p{27mm}p{18mm}Y}
\toprule
Stage & Upstream source & Split policy & Final inventory & Audit detail \\
\midrule
P1 & vetted VEH anchor contexts rendered through eight certification templates & three matched partitions under three seeds; subtype-balanced and not DOI-disjoint & 240 = 80/80/80 & legacy \texttt{dev}/\texttt{test\_id}/\texttt{test\_ood} labels are bookkeeping only; partitions use 20/19/19 anchor contexts with pairwise overlap 19/19/18 \\
P2 & 52 cleaned structure-review anchors from the extraction audit & optimizer-facing splits are DOI-disjoint: 26 oracle-domain anchors feed \texttt{dev}/\texttt{test\_id}, 26 projected out-of-domain-kept anchors feed \texttt{test\_ood} & 208 = 64/40/104 & 8 unit-suspect papers retained with explicit flags; admission report records source-anchor counts and no dropped task ids \\
P3 & dual-constraint trap projection from P2 & inherits P2 optimizer-facing split boundary & 156 = 64/40/52 & manifests record source P2 task bank, BKF table, split manifest, and trap policy \\
P4 & policy-conditioned ranking pools derived from P2 BKF records & inherits P2 optimizer-facing split boundary & 159 = 93/36/30 & 53 realized base pools across the release; manifests record profile counts, top-1 distinctness, and source-task lineage for each pool \\
\bottomrule
\end{tabularx}
\end{center}

The release package is designed around these manifests and is available through the anonymous artifact page at \href{https://anonymous.4open.science/r/diagbench-734D/README.md}{anonymous.4open.science/r/diagbench-734D}. It contains the exact task JSONL files and manifests for P1--P4, the oracle and evaluator code, the task-generation scripts, the admission and split reports, and the per-model JSONL run logs used for the main and appendix tables. The intended artifact boundary is therefore the same as the paper boundary: each reported benchmark table should be traceable to a specific manifest-backed snapshot rather than to an undocumented local run.

Contamination resistance in the current release is deliberately scoped. It means DOI/title-hash deduplication in the upstream paper registry, DOI-disjoint source-anchor separation for P2--P4, manifest-level lineage tracking, and evaluation on post-fix noleak P1 runs only. It does \emph{not} mean that we claim web-scale near-duplicate search over model pretraining corpora or model-specific pretraining decontamination. P1 is intentionally outside the source-disjoint claim and should be read as a matched robustness stress stage rather than as a held-out generalization benchmark.

\section{P3 Intervention Provenance}

The P3 intervention in the main paper is intentionally not a rerun over the entire 156-task bank. It is a frozen protocol audit on a 36-task subset selected by stratified sampling over \texttt{split} $\times$ \texttt{source\_p2\_subtype}. The selection rule is: draw three tasks per observed stratum, then top up to 12 tasks per split from the remaining pool while preserving split balance. This keeps every observed family visible while preventing the intervention table from collapsing to a single split.

\begin{table*}[t]
\caption{Full-bank versus intervention-subset composition for P3. The 36-task intervention subset is stratified rather than ad hoc: every observed split/subtype family is retained, and the final subset is exactly balanced across splits.}
\label{tab:appendix_p3_subset_provenance}
\centering
\small
\begin{tabular}{lp{50mm}p{50mm}}
\toprule
Split & Full 156-task bank & 36-task intervention subset \\
\midrule
\texttt{dev} & 64 tasks: boundary\_binding 20, paper\_like 17, power\_tight 12, resonance\_tuned 15 & 12 tasks: 3 / 3 / 3 / 3 across the same four strata \\
\texttt{test\_id} & 40 tasks: boundary\_binding 10, paper\_like 4, power\_tight 15, resonance\_tuned 11 & 12 tasks: 3 / 3 / 3 / 3 across the same four strata \\
\texttt{test\_ood} & 52 tasks: projected\_boundary 18, projected\_resonance 34 & 12 tasks: projected\_boundary 3, projected\_resonance 9 after split-balanced top-up \\
\bottomrule
\end{tabular}
\end{table*}

\section{Statistical Validity}

Headline gaps should be paired with uncertainty rather than interpreted as isolated point estimates. For continuous aggregates such as P1-Composite, P2 final feasible power ratio, and P4 Kendall-style scores, we use nonparametric bootstrap intervals over task rows. For binary proportions such as P3-Success, we report Wilson intervals. Paired bootstrap deltas are used for the main near ties and for the P3 intervention audit so that uncertainty reflects within-task covariance rather than independent resampling.

\begin{table*}[t]
\caption{Headline uncertainty table for the four main probes.}
\label{tab:appendix_ci_table}
\centering
\scriptsize
\setlength{\tabcolsep}{3.5pt}
\renewcommand{\arraystretch}{1.08}
\begin{tabular}{lp{23mm}p{18mm}p{23mm}p{25mm}p{36mm}}
\toprule
Probe & Metric & Best & 95\% CI & Key delta & Note \\
\midrule
P1 & P1-Composite & 0.574 & [0.500, 0.641] & qwen3-max $-$ gemini-3.1-pro-preview: +0.025 [-0.078, 0.121] & not a trivial propose-bias win \\
P2 & feasible power ratio & 0.3904 & [0.3598, 0.4235] & gemini-3.1-pro-preview $-$ claude-4.6-sonnet: +0.1510 [0.1117, 0.1891] & separates closure from anchoring \\
P3 & P3-Success & 42.9\% & [35.4, 50.8]\% & R1 $-$ gpt-5.4: +0.6 pts [-7.1, 8.3] & near tie at the recovery frontier \\
P4 & Kendall Tau & 0.887 & [0.862, 0.911] & gpt-5.4 $-$ qwen3.6-plus: +0.013 [-0.024, 0.050] & calibrates ranking claims \\
\bottomrule
\end{tabular}
\end{table*}

\begin{table*}[t]
\caption{Near-tie paired tests that matter most to the main paper.}
\label{tab:appendix_near_ties}
\centering
\scriptsize
\setlength{\tabcolsep}{3pt}
\renewcommand{\arraystretch}{1.08}
\begin{tabularx}{\textwidth}{p{32mm}p{38mm}Y}
\toprule
Pair & Test object & Why this pair matters \\
\midrule
deepseek-r1 vs gpt-5.4 & P3-Success / RecoveryQuality & current recovery near tie; mechanism differs even if win disappears \\
gpt-5.4 vs qwen3.6-plus & P4 Full Tau / Exact & structural ranking versus endpoint preference matching \\
qwen3-max vs gemini-3.1-pro-preview & P1-Composite / P2 headline & certification-versus-closure split \\
\bottomrule
\end{tabularx}
\end{table*}

\begin{table*}[t]
\caption{P3 intervention numeric results. Delta columns report paired task-bootstrap 95\% intervals for \texttt{state\_summary} versus raw history.}
\label{tab:appendix_p3_intervention_delta_ci}
\centering
\scriptsize
\setlength{\tabcolsep}{4pt}
\renewcommand{\arraystretch}{1.08}
\begin{tabular}{lcccc}
\toprule
Model & Succ. raw/sum. & $\Delta$ success & Casc. raw/sum. & $\Delta$ cascade \\
\midrule
deepseek-r1 & 55.6 / 58.3 & +2.8 pts [-11.1, +16.7] & 30.3 / 8.7 & -21.6 pts [-37.5, -7.2] \\
gpt-5.4 & 63.9 / 63.9 & +0.0 pts [-13.9, +13.9] & 29.4 / 0.0 & -29.4 pts [-45.5, -14.7] \\
gemini-3.1-pro-preview & 50.0 / 66.7 & +16.7 pts [-2.8, +36.1] & 46.9 / 25.0 & -21.9 pts [-43.6, +0.5] \\
claude-4.6-sonnet & 30.6 / 63.9 & +33.3 pts [+16.7, +50.0] & 5.7 / 4.3 & -1.4 pts [-8.6, +4.4] \\
\midrule
Model mean & 50.0 / 63.2 & +13.2 pts [+4.9, +21.5] & 28.1 / 9.5 & -18.6 pts [-26.6, -10.6] \\
\bottomrule
\end{tabular}
\end{table*}

\clearpage
\section{Metric Definitions}
\label{sec:metric_formulas}

This section defines the metrics used in the main paper and appendix. All scores are computed from persisted JSONL logs by deterministic evaluator code; no metric uses human annotation or LLM-as-a-judge.

\subsection{P1: Credible Triage}

P1 is a one-shot three-action decision problem with gold labels in
\(\{\texttt{propose\_design}, \texttt{declare\_infeasible}, \texttt{request\_missing\_info}\}\).

\textbf{Accuracy and Macro-F1.} Accuracy is micro-averaged over all P1 tasks. Macro-F1 is the unweighted mean of the three action-level F1 scores:
\begin{equation}
\text{Macro-F1}
= \frac{1}{3}\sum_{c\in\mathcal{C}} F_{1,c},
\qquad
F_{1,c}=
\frac{2\,\mathrm{Prec}_c\,\mathrm{Rec}_c}
{\mathrm{Prec}_c+\mathrm{Rec}_c}.
\end{equation}
If \(\mathrm{Prec}_c+\mathrm{Rec}_c=0\), the corresponding \(F_{1,c}\) is defined as zero.

\textbf{Discipline scores.} Let \(A_r\), \(M_r\), and \(I_r\) denote recall for \texttt{propose\_design}, \texttt{request\_missing\_info}, and \texttt{declare\_infeasible}. Let \(A_s\), \(M_s\), and \(I_s\) denote the corresponding spurious-action rates on tasks whose gold label is not that action. Then:
\begin{align}
\mathrm{ACS} &= A_r(1-A_s),\\
\mathrm{MDS} &= M_r(1-M_s),\\
\mathrm{IDS} &= I_r(1-I_s).
\end{align}
ACS penalizes habitual proposing, MDS penalizes over-requesting information, and IDS penalizes over-refusal.

\textbf{P1-Composite.} The P1 headline metric is the weighted certification score used by the evaluator:
\begin{equation}
\begin{aligned}
\mathrm{P1Comp}={}&
0.40\,F1_{\mathrm{3class}}
+0.20\,\mathrm{ACS}
+0.15\,\mathrm{MDS} \\
&+0.15\,\mathrm{IDS}
+0.10\,F1_{\mathrm{subtype}} .
\end{aligned}
\end{equation}

\subsection{P2: Constrained Repair / Search}

For a P2 trajectory, let \(\mathbf{x}_t\) be the design at step \(t\), \(V(\mathbf{x}_t)\ge0\) the total normalized oracle violation, and \(T\) the final evaluated step.

\textbf{First and final feasibility.} P2a is the fraction of tasks where \(V(\mathbf{x}_1)=0\). Final feasible rate is the fraction where \(V(\mathbf{x}_T)=0\).

\textbf{P2b final feasible power ratio.} For VEH tasks, the headline metric is:
\begin{equation}
\begin{aligned}
\mathrm{P2b}
= \frac{1}{N_{\mathrm{P2}}}\sum_{i=1}^{N_{\mathrm{P2}}}
&\mathbf{1}\!\left[V(\mathbf{x}^{(i)}_T)=0\right] \\
&\cdot
\frac{P_{\mathrm{out}}(\mathbf{x}^{(i)}_T)}
{P_{\mathrm{out}}(\mathbf{x}^{(i)}_{\mathrm{BKF}})} .
\end{aligned}
\end{equation}
The circuit audit uses the analogous final feasible objective score returned by the circuit oracle.

\textbf{Conditional objective ratio.}
\begin{equation}
\mathrm{CondObj} =
\frac{
\sum_{i:\,V(\mathbf{x}^{(i)}_T)=0}
\mathrm{obj}(\mathbf{x}^{(i)}_T)}
{|\{i: V(\mathbf{x}^{(i)}_T)=0\}|}.
\end{equation}

\textbf{Trajectory diagnostics.} Violation reduction consistency is the fraction of transitions with \(V(\mathbf{x}_{t+1})\le V(\mathbf{x}_t)-10^{-6}\). Directed repair rate is the fraction of violated-constraint updates that move the relevant oracle metric in the feedback-implied direction. Feasibility preservation is the fraction of feasible-to-feasible transitions. Mean log edit delta is:
\begin{equation}
\delta_{\mathrm{edit}}
=
\frac{1}{|\mathcal{D}|}\sum_{d\in\mathcal{D}}
\frac{\left|\log(x_{t+1,d}/x_{t,d})\right|}
{\log(x_d^{\max}/x_d^{\min})}.
\end{equation}
Over-edit and no-op rates threshold this quantity in the corresponding evaluator.

\subsection{P3: Post-Trap Recovery}

P3 starts from a corrupted trajectory state. Escape rate is the fraction of tasks where at least one step reduces violation relative to the trapped state. Cascade rate is the fraction of escaped tasks where a new structurally coupled violation appears after the escape move. Dead-budget rate is the fraction of tasks with no meaningful post-trap design change.

\textbf{P3-Success and recovery quality.}
\begin{equation}
\mathrm{P3Success}
= \frac{1}{N_{\mathrm{P3}}}
\sum_i \mathbf{1}\!\left[V(\mathbf{x}^{(i)}_T)=0\right].
\end{equation}
\begin{equation}
\mathrm{RecoveryQuality}
=
\frac{
\sum_{i:\,V(\mathbf{x}^{(i)}_T)=0}
\mathrm{obj}(\mathbf{x}^{(i)}_T)}
{|\{i: V(\mathbf{x}^{(i)}_T)=0\}|}.
\end{equation}
The raw-history versus state-summary delta is the paired task-level difference in P3-Success between the two prompt representations, reported with paired bootstrap intervals.

\subsection{P4: Policy-Conditioned Ranking}

Each P4 task contains five oracle-feasible candidates. Full Kendall \(\tau\) is Kendall \(\tau_b\) between the model ranking and oracle ranking. Exact match, top-1 accuracy, top-2 set accuracy, pairwise accuracy, and policy-flip accuracy are computed directly from the predicted order.

\textbf{BARS.} For the VEH P4-full bank, the headline BARS score is computed only on balanced-active rows:
\begin{equation}
\begin{aligned}
\mathrm{BARS}
= {}&0.55\,\tau_{\mathrm{BA}}
+0.25\,\mathrm{PairAcc}_{\mathrm{BA}} \\
&+0.20\,\mathrm{Exact}_{\mathrm{BA}},
\end{aligned}
\end{equation}
where BA denotes the balanced-active subset. For the circuit audit, the analogous ranking score uses scaled Kendall \((\tau+1)/2\), policy-flip accuracy, and exact match with the same weights.

\subsection{Response-Control Profile Scores}

Profile scores are log-derived diagnostics. For each dimension \(D\), the evaluator collects a fixed indicator set \(\mathcal{M}_D\), orients every indicator so higher is better, min-max scales it across the compared models, and averages:
\begin{equation}
\mathrm{Profile}_D(m)
=
\frac{1}{|\mathcal{M}_D|}
\sum_{q\in\mathcal{M}_D} z_q(m).
\end{equation}

\begin{center}
\captionof{table}{Indicators used for response-control profile extraction. All indicators are direction-normalized before averaging.}
\label{tab:appendix_profile_indicator_sets}
\scriptsize
\setlength{\tabcolsep}{4pt}
\renewcommand{\arraystretch}{1.08}
\begin{tabular}{lp{115mm}}
\toprule
Dimension & Indicators \\
\midrule
Action prior & P1 action-distribution alignment, macro-F1, missing recall, infeasible recall, propose precision, non-invalid rate \\
Edit style & P2 bounded-local-edit rate, feasibility preservation, directed repair, final feasible rate, non-destructive edit rate, protocol-valid rate \\
Feedback obedience & P2 violation reduction, P2 utility improvement, P2 best-so-far AUC, P3 violation reduction, P3 post-feedback feasible rate \\
State trust & P3 trap escape, explicit replan, escape quality, non-cascade rate, non-dead-budget rate, state-summary success gain, state-summary cascade reduction \\
Preference execution & P4 scaled full Tau, balanced-active BARS, balanced-active policy-sensitive pair accuracy, exact match, top-1 accuracy, non-Pareto-violation rate, non-parse-error rate \\
\bottomrule
\end{tabular}
\end{center}


\section{Calibration References}

Weak baselines are useful here because the benchmark is deliberately structured. For P1, the main concern is degenerate label-prior matching. For P2, the main concern is whether headline gains reflect true improvement beyond the literature-derived seed or only feasibility preservation. For P3, the relevant weak behavior is continuity with the corrupted state rather than explicit reset. For P4, the relevant floor is chance agreement with a five-item ordering.

\begin{table*}[t]
\caption{Calibration references corresponding to the main-text weak baselines. These are not competitors to the frontier roster; they anchor how much of each stage can be explained by trivial behavior.}
\label{tab:appendix_calibration}
\centering
\small
\begin{tabular}{lp{30mm}p{23mm}p{60mm}}
\toprule
Probe & Baseline reference & Score & Interpretation \\
\midrule
P1 & always-\texttt{propose\_design} on the dev split & 56.3\% accuracy, 0.240 macro-F1, 0.0 worst-action recall & a degenerate action prior can look superficially acceptable on raw accuracy, which is why P1-Composite and subtype-balanced reporting are necessary \\
P2 & anchor-fixed heuristic on the 208-task main bank & 100\% feasible coverage, 0.3065 mean power ratio & many frontier models fail to preserve even the trivial feasible anchor, while gemini-3.1-pro-preview is the only one to beat this heuristic on the headline utility ratio \\
P3 & last-state continuation / no-reset heuristic & 0\% explicit reset by construction; weak reference for escape, cascade, and recovery columns & the useful floor is not a single scalar success rate but the failure mode produced by trusting corrupted state; P3 therefore reports escape, cascade, dead-budget, and recovery jointly \\
P4 & uniform random permutation over five feasible candidates & $\mathbb{E}[\tau]=0$, exact 0.83\%, top-1 20\% & observed ranking scores are far above chance, so P4 is not a random-agreement artifact \\
\bottomrule
\end{tabular}
\end{table*}

\clearpage
\section{Full Leaderboard Tables}

The appendix carries the full all-model tables for each main probe so that the main-paper mechanism claims remain auditable. All four tables use the frozen 12-model complete roster; runs without complete P1--P4 coverage are excluded rather than imputed.

This block should be read as a diagnostic audit rather than four extra leaderboards. Read each table column-wise: the headline metric gives the rank, while the diagnostic columns explain why that rank occurred and what failure mode is being exposed. The formulas are collected in Appendix~\ref{sec:metric_formulas}; the paragraphs below summarize how each metric is computed, what behavior it is intended to reveal, and what conclusion follows from the observed table.

\begin{table*}[t]
\caption{Interpretive map for the appendix full-table block. Each table is tied to a benchmark claim rather than presented as metric inventory alone.}
\label{tab:appendix_fulltable_map}
\centering
\small
\begin{tabular}{lp{21mm}p{16mm}p{16mm}p{16mm}p{16mm}p{18mm}}
\toprule
Probe & Core columns & Diagnostic 1 & Diagnostic 2 & Diagnostic 3 & Diagnostic 4 & Primary mechanism claim served \\
\midrule
P1 & Acc / Macro-F1 / ACS / MDS / IDS / subtype F1 / composite & over-act prior & hard infeasibility misses & missing-info discipline & balanced selectivity & raw accuracy does not equal credible triage \\
P2 & P2a / P2b / final feasible rate / conditional ratio / Impr. (P2c) / queries & anchoring & closure & best-so-far improvement & query efficiency & start is not finish \\
P3 & success / recovery quality / trap escape / cascade / dead budget / explicit replan & no-escape failure & escaped-but-unrecovered & false recovery & post-escape quality & escape is not recovery \\
P4 & Full Tau / Exact / Top-1 / Top-2 / Pareto Tau / BARS & structural order & strict endpoint ranking & policy-active ranking & parse stability & structural ranking is not policy sensitivity \\
\bottomrule
\end{tabular}
\end{table*}

\Needspace{14\baselineskip}
\paragraph{P1 metrics and readout.}
P1 is a one-shot triage task. Accuracy counts exact agreement with the oracle label; macro-F1 averages the three action labels so that a model cannot score well by favoring the majority action. ACS, MDS, and IDS are discipline scores for \texttt{propose\_design}, \texttt{request\_missing\_info}, and \texttt{declare\_infeasible}; each multiplies recall by one minus the corresponding spurious-action rate. Subtype F1 checks whether performance holds across near-feasible, missing-info, and infeasible subtypes. The headline P1-Composite combines these terms, so it rewards balanced entry discipline rather than raw willingness to produce a design. The table shows why this is necessary: several models have similar raw accuracy, but qwen3-max and gemini-3.1-pro-preview separate because their discipline scores are more balanced, while claude-4.6-sonnet exposes over-refusal or action-mismatch behavior that accuracy alone would under-explain.

\begin{center}
\captionof{table}{P1 full table. Compact headers keep certification metrics readable.}
\label{tab:appendix_p1_full}
\scriptsize
\setlength{\tabcolsep}{3.5pt}
\renewcommand{\arraystretch}{1.08}
\begin{tabular}{llllllll}
\toprule
Model & Acc & F1 & ACS & MDS & IDS & Subtype F1 & Comp. \\
\midrule
qwen3-max & 65.4\% & 0.647 & 0.429 & 0.533 & 0.486 & 0.768 & 0.574 \\
gemini-3.1-pro-preview & 69.2\% & 0.636 & 0.430 & 0.667 & 0.218 & 0.762 & 0.549 \\
\texttt{o4-mini} & 70.0\% & 0.623 & 0.388 & 0.396 & 0.361 & 0.776 & 0.518 \\
deepseek-r1 & 66.2\% & 0.594 & 0.360 & 0.510 & 0.298 & 0.730 & 0.504 \\
gpt-5.4 & 43.3\% & 0.486 & 0.224 & 0.657 & 0.223 & 0.569 & 0.428 \\
hunyuan-hy3-preview & 66.2\% & 0.539 & 0.295 & 0.333 & 0.200 & 0.700 & 0.425 \\
deepseek-v3 & 64.6\% & 0.483 & 0.206 & 0.352 & 0.100 & 0.668 & 0.369 \\
\texttt{llama-3.3-70b} & 62.5\% & 0.467 & 0.242 & 0.200 & 0.192 & 0.669 & 0.361 \\
mimo-v2.5-pro & 61.7\% & 0.433 & 0.152 & 0.244 & 0.083 & 0.644 & 0.317 \\
\texttt{qwen3.6-plus} & 61.3\% & 0.422 & 0.143 & 0.200 & 0.100 & 0.641 & 0.306 \\
deepseek-v4-pro & 61.3\% & 0.422 & 0.143 & 0.200 & 0.100 & 0.641 & 0.306 \\
claude-4.6-sonnet & 57.5\% & 0.305 & 0.057 & 0.000 & 0.100 & 0.583 & 0.207 \\
\bottomrule
\end{tabular}
\end{center}

\noindent\textbf{Audit note.} The exact tie between \texttt{qwen3.6-plus} and deepseek-v4-pro is preserved from the frozen scoring snapshot rather than broken by post-hoc rounding or reranking. Exact vector ties are therefore treated as tied diagnostic rows, not as evidence for an ordering between the two models.

\noindent\textbf{Takeaway.} P1 separates credible entry control from raw action frequency: the strongest rows combine reasonable proposing with missing-information and infeasibility discipline, while weak rows expose over-refusal or over-action even when accuracy is not catastrophic.

\Needspace{14\baselineskip}
\paragraph{P2 metrics and readout.}
P2 is an iterative repair task under verifier feedback. P2a is first-step feasibility, which measures whether the first edit lands directly in a feasible region. P2b is the headline final feasible power ratio: the final design contributes only if it is feasible, and its output power is normalized by the benchmark-known feasible reference. Final feasibility separates closure from objective quality, conditional ratio reports quality only among successful repairs, improvement measures best-so-far utility gain, AUC summarizes trajectory quality across the budget, and queries records how many oracle calls were needed. The main conclusion is that search closure and utility are not the same: gemini-3.1-pro-preview dominates P2b because it combines high final feasibility with high feasible utility, while models with respectable first-step feasibility can still finish with low final ratio if their subsequent edits overfit one constraint or lose power.

\begin{center}
\captionof{table}{P2 full table. P2a is first-step feasibility; P2b is the headline final feasible power ratio.}
\label{tab:appendix_p2_full}
\scriptsize
\setlength{\tabcolsep}{3.5pt}
\renewcommand{\arraystretch}{1.08}
\begin{tabular}{llllllll}
\toprule
Model & P2a & P2b & Final feas. & Cond. ratio & Impr. & AUC & Queries \\
\midrule
gemini-3.1-pro-preview & 26.0\% & 0.3904 & 96.2\% & 0.4101 & 0.7902 & 0.3388 & 2.63 \\
deepseek-r1 & 29.8\% & 0.2413 & 67.8\% & 0.3585 & 0.7901 & 0.2999 & 3.47 \\
claude-4.6-sonnet & 13.0\% & 0.2394 & 62.5\% & 0.3891 & 0.7037 & 0.2797 & 4.41 \\
\texttt{qwen3.6-plus} & 30.8\% & 0.2063 & 61.1\% & 0.3432 & 0.6466 & 0.2880 & 3.87 \\
qwen3-max & 23.6\% & 0.1955 & 55.8\% & 0.3537 & 0.6649 & 0.2894 & 4.21 \\
hunyuan-hy3-preview & 2.4\% & 0.1609 & 48.6\% & 0.3380 & 0.8033 & 0.2480 & 4.79 \\
deepseek-v3 & 7.2\% & 0.1565 & 44.2\% & 0.3576 & 0.5907 & 0.2766 & 4.68 \\
\texttt{o4-mini} & 22.6\% & 0.1551 & 42.8\% & 0.3665 & 0.7571 & 0.3119 & 3.59 \\
gpt-5.4 & 7.2\% & 0.1329 & 36.1\% & 0.3841 & 0.6828 & 0.2708 & 4.75 \\
deepseek-v4-pro & 30.3\% & 0.1294 & 41.3\% & 0.3205 & 0.6167 & 0.2935 & 4.26 \\
\texttt{llama-3.3-70b} & 23.6\% & 0.1197 & 37.0\% & 0.3277 & 0.6673 & 0.2895 & 4.38 \\
mimo-v2.5-pro & 26.4\% & 0.1073 & 33.2\% & 0.3282 & 0.6714 & 0.3075 & 4.50 \\
\bottomrule
\end{tabular}
\end{center}

\noindent\textbf{Takeaway.} P2 is won by feasible closure plus useful final power, not by first-step feasibility alone; gemini-3.1-pro-preview's lead is therefore a repair-trajectory result rather than an anchoring artifact.

\Needspace{14\baselineskip}
\paragraph{P3 metrics and readout.}
P3 starts from a corrupted trajectory. Success is final feasibility at the end of the recovery budget; recovery quality reports the objective quality among successful recoveries; first feasible identifies whether a model ever reaches a feasible state early in the recovery trace. Escape measures whether the model moves away from the known trap, cascade records whether the escape introduces a new coupled violation, dead-budget marks trajectories with no meaningful recovery action, and replan counts explicit reset or replanning behavior. These columns show that P3 is not ordinary repair: claude-4.6-sonnet is the clearest diagnostic row, with 96.8\% escape but only 0.6\% first-feasible and 16.0\% final success. It often notices that the raw history is wrong, but it rarely stabilizes back into a feasible design. deepseek-r1 and gpt-5.4 combine high escape with better final recovery; hunyuan-hy3-preview and mimo-v2.5-pro lead success through different trade-offs in recovery quality and cascade control.

\begin{center}
\captionof{table}{P3 full table. Columns separate recovery success, escape, cascade, dead-budget, and replanning behavior.}
\label{tab:appendix_p3_full}
\scriptsize
\setlength{\tabcolsep}{3.5pt}
\renewcommand{\arraystretch}{1.08}
\begin{tabular}{llllllll}
\toprule
Model & Succ. & RecQ & First feas. & Escape & Cascade & Dead & Replan \\
\midrule
hunyuan-hy3-preview & 47.4\% & 0.4876 & 40.4\% & 77.6\% & 9.1\% & 25.3\% & 7.7\% \\
mimo-v2.5-pro & 45.5\% & 0.6263 & 41.0\% & 63.5\% & 4.0\% & 26.7\% & 17.9\% \\
deepseek-r1 & 42.9\% & 0.2708 & 22.4\% & 92.3\% & 26.4\% & 25.1\% & 26.3\% \\
gpt-5.4 & 42.3\% & 0.3723 & 26.3\% & 90.4\% & 21.3\% & 25.2\% & 6.4\% \\
gemini-3.1-pro-preview & 37.2\% & 0.3194 & 26.9\% & 92.3\% & 50.7\% & 25.8\% & 10.9\% \\
deepseek-v3 & 35.3\% & 0.2868 & 23.1\% & 82.7\% & 27.1\% & 27.6\% & 57.7\% \\
qwen3-max & 30.1\% & 0.2612 & 22.4\% & 85.9\% & 12.7\% & 25.7\% & 49.4\% \\
deepseek-v4-pro & 28.8\% & 0.1986 & 27.6\% & 90.4\% & 17.0\% & 24.3\% & 99.4\% \\
\texttt{qwen3.6-plus} & 27.6\% & 0.1453 & 12.8\% & 94.9\% & 23.0\% & 25.6\% & 0.6\% \\
\texttt{o4-mini} & 26.3\% & 0.1389 & 12.2\% & 69.2\% & 17.6\% & 26.6\% & 8.3\% \\
claude-4.6-sonnet & 16.0\% & 0.0960 & 0.6\% & 96.8\% & 5.3\% & 25.7\% & 0.0\% \\
\texttt{llama-3.3-70b} & 3.2\% & 0.0926 & 2.6\% & 34.6\% & 3.7\% & 14.3\% & 0.0\% \\
\bottomrule
\end{tabular}
\end{center}

\noindent\textbf{Takeaway.} P3 shows that escaping a corrupted trajectory is not the same as recovering: high escape with low success marks continuity or cascade failure, while successful recovery requires both reset and stabilization.

\Needspace{14\baselineskip}
\paragraph{P4 metrics and readout.}
P4 removes search from the problem: all candidates are oracle-feasible before the model ranks them. Full Kendall Tau measures global order agreement with the oracle policy ranking; exact match requires the entire five-item order to be correct; top-1 and top-2 set accuracy measure whether the highest-priority choices are preserved; Pareto Tau isolates ranking consistency on Pareto-relevant comparisons; BARS emphasizes balanced-active policy-sensitive rows; parse and violation rates record formatting failures and attempts to alter fixed candidates. The table shows that P4 is not reducible to one ranking statistic. gpt-5.4 leads Full Tau, qwen3.6-plus is strongest on exact match, and gemini-3.1-pro-preview has a high BARS despite lower Tau, indicating that policy execution has multiple sub-behaviors rather than a single scalar notion of ``ranking skill.''

\begin{center}
\captionof{table}{P4 full table. Compact notation keeps ranking, Pareto, and formatting diagnostics in one table.}
\label{tab:appendix_p4_full}
\scriptsize
\setlength{\tabcolsep}{3.5pt}
\renewcommand{\arraystretch}{1.08}
\begin{tabular}{llllllll}
\toprule
Model & Tau & Exact & Top1 & Top2 & Pareto & BARS & Parse / viol. \\
\midrule
gpt-5.4 & 0.887 & 57.2\% & 83.0\% & 81.1\% & 1.000 & 0.6258 & 0.0\% / 0.0\% \\
\texttt{qwen3.6-plus} & 0.877 & 66.3\% & 79.1\% & 83.4\% & 1.000 & 0.6912 & 0.0\% / 0.0\% \\
deepseek-v3 & 0.860 & 54.7\% & 77.4\% & 79.2\% & 0.984 & 0.5751 & 0.0\% / 0.8\% \\
mimo-v2.5-pro & 0.843 & 52.8\% & 72.3\% & 79.2\% & 0.993 & 0.6012 & 0.0\% / 0.3\% \\
claude-4.6-sonnet & 0.840 & 56.6\% & 72.3\% & 79.2\% & 1.000 & 0.6781 & 0.0\% / 0.0\% \\
hunyuan-hy3-preview & 0.839 & 50.9\% & 77.4\% & 79.9\% & 0.943 & 0.5610 & 0.0\% / 2.8\% \\
qwen3-max & 0.835 & 49.7\% & 78.0\% & 76.7\% & 0.983 & 0.5367 & 0.0\% / 0.8\% \\
deepseek-r1 & 0.833 & 56.0\% & 71.7\% & 76.7\% & 1.000 & 0.6835 & 0.0\% / 0.0\% \\
gemini-3.1-pro-preview & 0.824 & 54.1\% & 71.7\% & 78.6\% & 0.975 & 0.6988 & 1.3\% / 1.3\% \\
deepseek-v4-pro & 0.794 & 43.4\% & 60.4\% & 74.2\% & 0.977 & 0.5832 & 0.0\% / 1.2\% \\
\texttt{o4-mini} & 0.780 & 50.0\% & 60.0\% & 72.0\% & 1.000 & 0.5988 & 0.0\% / 0.0\% \\
\texttt{llama-3.3-70b} & 0.714 & 34.0\% & 54.1\% & 65.4\% & 0.959 & 0.5386 & 0.0\% / 2.0\% \\
\bottomrule
\end{tabular}
\end{center}

\noindent\textbf{Takeaway.} P4 ranking has multiple sub-behaviors: global order quality, exact endpoint order, top-choice preservation, and policy-sensitive BARS can diverge, which is why P4 is not reducible to one scalar preference score.

\section{Split-Resolved Tables}

This block checks that the main-paper claims are not driven by a single split. The P2 closure pattern and P4 non-monotonic ranking pattern remain visible across splits, while P3 contains genuine near ties that should be read with the confidence intervals in Appendix~\ref{tab:appendix_ci_table}. P1 uses matched partitions rather than DOI-disjoint source splits, so its split labels audit robustness rather than held-out generalization. The split-resolved and stage-gap tables use the frozen common-coverage split snapshot; later-added full-bank rows for hunyuan-hy3-preview, mimo-v2.5-pro, and deepseek-v4-pro are included in the main full tables but not imputed into split/stage-gap views. Exact split-invariant rows, such as \texttt{qwen3.6-plus} and claude-4.6-sonnet in P1, are retained as audit signals from the released snapshot. Each table reports the main stage headline plus one supporting diagnostic in a split-resolved view.

\begin{table*}[t]
\caption{Split and subtype reporting plan. The appendix presents split-resolved views as mechanism audits rather than as extra leaderboards.}
\label{tab:appendix_split_subtype_map}
\centering
\small
\begin{tabular}{lp{28mm}p{32mm}p{42mm}}
\toprule
Probe & Split-resolved view & Subtype families & Main claim audited \\
\midrule
P1 & matched-partition composite and diagnostics (legacy \texttt{dev}/\texttt{test\_id}/\texttt{test\_ood} labels) & infeasible-margin, infeasible-structural, missing-info families & certification robustness is not driven by one matched partition or blocker family \\
P2 & dev / test\_id / test\_ood headline and feasible rate & paper-like, resonance-tuned, power-tight, boundary-binding & closure dominates anchoring as the main search distinction \\
P3 & dev / test\_id / test\_ood success and recovery quality & trap families and cascade diagnostics & corrupted-state recovery is not reducible to nominal repair quality \\
P4 & dev / test\_id / test\_ood Full Tau and Exact & balanced, performance-first, reliability-first, BARS & structural ranking and policy-sensitive ranking remain partially separable \\
\bottomrule
\end{tabular}
\end{table*}

\begin{table*}[t]
\caption{P1 matched-partition table. Each cell reports P1-Composite and raw accuracy.}
\label{tab:appendix_p1_split}
\centering
\scriptsize
\setlength{\tabcolsep}{4pt}
\renewcommand{\arraystretch}{1.08}
\begin{tabular}{lccc}
\toprule
Model & dev & test\_id & test\_ood \\
\midrule
qwen3-max & 0.532 / 61.3\% & 0.622 / 70.0\% & 0.567 / 65.0\% \\
gemini-3.1-pro-preview & 0.545 / 70.0\% & 0.525 / 66.2\% & 0.578 / 71.2\% \\
\texttt{o4-mini} & 0.542 / 71.2\% & 0.509 / 70.0\% & 0.502 / 68.8\% \\
deepseek-r1 & 0.517 / 68.8\% & 0.512 / 66.2\% & 0.482 / 63.7\% \\
gpt-5.4 & 0.374 / 37.5\% & 0.469 / 47.5\% & 0.440 / 45.0\% \\
deepseek-v3 & 0.357 / 63.7\% & 0.385 / 65.0\% & 0.364 / 65.0\% \\
\texttt{llama-3.3-70b} & 0.362 / 61.3\% & 0.393 / 65.0\% & 0.328 / 61.3\% \\
\texttt{qwen3.6-plus} & 0.306 / 61.3\% & 0.306 / 61.3\% & 0.306 / 61.3\% \\
claude-4.6-sonnet & 0.207 / 57.5\% & 0.207 / 57.5\% & 0.207 / 57.5\% \\
\bottomrule
\end{tabular}
\end{table*}

\begin{table*}[t]
\caption{P2 split-resolved table. Each cell reports the final feasible power ratio and final feasible coverage.}
\label{tab:appendix_p2_split}
\centering
\scriptsize
\setlength{\tabcolsep}{4pt}
\renewcommand{\arraystretch}{1.08}
\begin{tabular}{lccc}
\toprule
Model & dev & test\_id & test\_ood \\
\midrule
gemini-3.1-pro-preview & 0.4017 / 98.4\% & 0.4324 / 90.0\% & 0.3672 / 97.1\% \\
deepseek-r1 & 0.2560 / 62.5\% & 0.3049 / 82.5\% & 0.2078 / 65.4\% \\
claude-4.6-sonnet & 0.2199 / 54.7\% & 0.3095 / 75.0\% & 0.2245 / 62.5\% \\
\texttt{qwen3.6-plus} & 0.2355 / 57.8\% & 0.2316 / 65.0\% & 0.1786 / 61.5\% \\
qwen3-max & 0.2182 / 50.0\% & 0.2486 / 60.0\% & 0.1612 / 57.7\% \\
deepseek-v3 & 0.2052 / 48.4\% & 0.1761 / 47.5\% & 0.1189 / 40.4\% \\
\texttt{o4-mini} & 0.1587 / 42.2\% & 0.1220 / 35.0\% & 0.1655 / 46.2\% \\
gpt-5.4 & 0.0991 / 26.6\% & 0.0541 / 12.5\% & 0.1841 / 51.0\% \\
\texttt{llama-3.3-70b} & 0.1781 / 43.8\% & 0.1046 / 32.5\% & 0.0896 / 34.6\% \\
\bottomrule
\end{tabular}
\end{table*}

\begin{table*}[t]
\caption{P3 split-resolved table. Each cell reports final recovered feasible rate and post-escape recovery quality.}
\label{tab:appendix_p3_split}
\centering
\scriptsize
\setlength{\tabcolsep}{4pt}
\renewcommand{\arraystretch}{1.08}
\begin{tabular}{lccc}
\toprule
Model & dev & test\_id & test\_ood \\
\midrule
deepseek-r1 & 45.3\% / 0.3103 & 37.5\% / 0.2105 & 44.2\% / 0.2708 \\
gpt-5.4 & 48.4\% / 0.4322 & 40.0\% / 0.3816 & 36.5\% / 0.2841 \\
gemini-3.1-pro-preview & 43.8\% / 0.3770 & 32.5\% / 0.2917 & 32.7\% / 0.2660 \\
deepseek-v3 & 42.2\% / 0.2368 & 27.5\% / 0.3571 & 32.7\% / 0.3068 \\
qwen3-max & 39.1\% / 0.3182 & 22.5\% / 0.1912 & 25.0\% / 0.2444 \\
\texttt{qwen3.6-plus} & 28.1\% / 0.1774 & 27.5\% / 0.1389 & 26.9\% / 0.1100 \\
\texttt{o4-mini} & 28.1\% / 0.1556 & 30.0\% / 0.1667 & 21.2\% / 0.0972 \\
claude-4.6-sonnet & 20.3\% / 0.1270 & 10.0\% / 0.0513 & 15.4\% / 0.0918 \\
\texttt{llama-3.3-70b} & 4.7\% / 0.1429 & 0.0\% / 0.0000 & 3.8\% / 0.1429 \\
\bottomrule
\end{tabular}
\end{table*}

\begin{table*}[t]
\caption{P4 split-resolved table. Each cell reports Full Tau and Exact.}
\label{tab:appendix_p4_split}
\centering
\scriptsize
\setlength{\tabcolsep}{4pt}
\renewcommand{\arraystretch}{1.08}
\begin{tabular}{lccc}
\toprule
Model & dev & test\_id & test\_ood \\
\midrule
gpt-5.4 & 0.886 / 53.8\% & 0.900 / 58.3\% & 0.873 / 66.7\% \\
\texttt{qwen3.6-plus} & 0.867 / 65.6\% & 0.870 / 57.5\% & 0.920 / 80.0\% \\
deepseek-v3 & 0.888 / 57.0\% & 0.783 / 47.2\% & 0.867 / 56.7\% \\
claude-4.6-sonnet & 0.832 / 59.1\% & 0.833 / 41.7\% & 0.873 / 66.7\% \\
qwen3-max & 0.832 / 49.5\% & 0.850 / 47.2\% & 0.827 / 53.3\% \\
deepseek-r1 & 0.834 / 57.0\% & 0.800 / 44.4\% & 0.867 / 66.7\% \\
gemini-3.1-pro-preview & 0.826 / 55.9\% & 0.778 / 38.9\% & 0.873 / 66.7\% \\
\texttt{o4-mini} & 0.857 / 57.1\% & 0.750 / 44.4\% & 0.857 / 71.4\% \\
\texttt{llama-3.3-70b} & 0.720 / 35.5\% & 0.667 / 22.2\% & 0.753 / 43.3\% \\
\bottomrule
\end{tabular}
\end{table*}

\section{Reasoning Coverage and Model Style}

This block remains secondary evidence rather than a substitute for the model-by-model tables. It is a negative control for the response-control profile account: if the paper were only measuring whether thinking mode was on, the stage patterns should collapse under this grouping. Instead, thinking mode helps some workflow regimes more than others and still does not erase stage dissociation.

\begin{table*}[t]
\caption{Reasoning/thinking run-mode audit for the 12 complete P1--P4 model runs. The column records whether thinking mode was used in the reported run.}
\label{tab:appendix_reasoning_grouping}
\centering
\scriptsize
\setlength{\tabcolsep}{3.5pt}
\renewcommand{\arraystretch}{1.08}
\begin{tabular}{lp{33mm}p{17mm}p{54mm}}
\toprule
Model & Public / run evidence & Think used & Caveat \\
\midrule
\texttt{o4-mini} & OpenAI documents it as a reasoning model~\citep{openai2026o4mini} & Yes & reasoning-labeled row \\
gemini-3.1-pro-preview & Google documents internal thinking / thought-summary support~\citep{google2026geminithinking} & Yes & thinking support active in the reported run \\
deepseek-r1 & DeepSeek documentation positions R1 as a reasoning-oriented model~\citep{deepseek2025r1} & Yes & key within-family contrast is R1 versus V3 \\
hunyuan-hy3-preview & public provider pages describe Hy3 as a fast/slow or multi-mode reasoning model~\citep{siliconflow2026hy3preview} & Yes & run manifests are retained as the source of truth for provider settings \\
deepseek-v4-pro & DeepSeek documentation describes V4 as supporting Thinking and Non-Thinking modes~\citep{deepseek2026v4preview} & No & manifest records \texttt{thinking=disabled} for JSON/action-schema stability \\
mimo-v2.5-pro & Xiaomi Mimo API row treated as a thinking-capable provider model in the run plan & No & manifest records \texttt{thinking=disabled} for JSON/action-schema stability \\
qwen3-max & Alibaba documents Qwen3-Max as hybrid-thinking with thinking disabled by default~\citep{alibaba2026qwenDeepThinking} & No & DashScope call used provider default; no explicit \texttt{enable\_thinking=True} override \\
\texttt{qwen3.6-plus} & Qwen documentation describes explicit \texttt{enable\_thinking=True} usage~\citep{qwen2026qwen36plus} & No & DashScope call used provider default without an explicit \texttt{enable\_thinking} flag \\
gpt-5.4 & frontier GPT row with provider-default reasoning effort in our logs & No & reported as no explicit thinking-mode run \\
deepseek-v3 & no reasoning-first grouping used here & No & used as the within-family non-R1 contrast \\
claude-4.6-sonnet & no reasoning-first grouping used here & No & grouped by observed profile, not by hidden implementation details \\
\texttt{llama-3.3-70b} & no reasoning-first grouping used here & No & open model baseline row in the complete P1--P4 roster \\
\bottomrule
\end{tabular}
\end{table*}

\begin{table}[t]
\caption{Stage-wise group means on the 12 complete model runs, grouped only by whether thinking mode was used in the reported run.}
\label{tab:appendix_reasoning_groupmeans}
\centering
\scriptsize
\setlength{\tabcolsep}{5pt}
\renewcommand{\arraystretch}{1.08}
\begin{tabular}{lcc}
\toprule
Stage & Think & No-think \\
\midrule
P1 Composite & 0.499 & 0.359 \\
P2 headline & 0.237 & 0.161 \\
P3 Success (\%) & 38.5 & 28.6 \\
P4 Full Tau & 0.819 & 0.831 \\
\bottomrule
\end{tabular}
\end{table}

% Auto-generated by scripts/quantify_response_control_profiles.py
% Columns: Model, Action prior, Edit style, Feedback obedience, State trust, Preference execution
\begin{table*}[t]
\caption{Stable response-control profile scores on the 12 complete P1--P4 runs. Scores are diagnostic boundary-fit indicators, not a monolithic leaderboard.}
\label{tab:profile_quantification}
\centering
\scriptsize
\setlength{\tabcolsep}{4pt}
\renewcommand{\arraystretch}{1.08}
\begin{tabular}{lccccc}
\toprule
Model & Action & Edit & Feedback & State & Preference \\
\midrule
model_A & 0.737 & 0.571 & 0.407 & 0.646 & 0.679 \\
model_B & 0.668 & 0.753 & 0.508 & 0.425 & 0.729 \\
model_C & 0.657 & 0.482 & 0.406 & 0.453 & 0.743 \\
model_D & 0.652 & 0.667 & 0.455 & 0.455 & 0.728 \\
model_F & 0.635 & 0.470 & 0.391 & 0.495 & 0.666 \\
model_M & 0.564 & 0.485 & 0.429 & 0.599 & 0.688 \\
model_H & 0.536 & 0.493 & 0.389 & 0.629 & 0.708 \\
Llama-3.3 & 0.528 & 0.586 & 0.326 & 0.452 & 0.622 \\
model_L & 0.499 & 0.385 & 0.391 & 0.627 & 0.707 \\
model_K & 0.491 & 0.404 & 0.332 & 0.737 & 0.666 \\
model_G & 0.491 & 0.613 & 0.404 & 0.523 & 0.761 \\
model_E & 0.428 & 0.622 & 0.392 & 0.443 & 0.730 \\
\bottomrule
\end{tabular}
\end{table*}


\begin{table*}[t]
\caption{Paired model-style contrasts.}
\label{tab:appendix_modelstyle_pairs}
\centering
\scriptsize
\setlength{\tabcolsep}{4pt}
\renewcommand{\arraystretch}{1.08}
\begin{tabular}{lp{27mm}p{27mm}p{47mm}}
\toprule
Contrast & First stronger & Second stronger & Isolated contrast \\
\midrule
deepseek-r1 vs deepseek-v3 & P1 / P2 / P3 & P4 & within-family evidence that reasoning-style gains cluster in certification and closure rather than ranking \\
gemini-3.1-pro-preview vs gpt-5.4 & P2 & P4 & closure-versus-ranking split across two frontier systems \\
\texttt{o4-mini} vs gpt-5.4 & P1 / weakly on P2 & P3 / P4 & explicit reasoning coverage does not automatically dominate recovery or ranking \\
\bottomrule
\end{tabular}
\end{table*}

\begin{table*}[t]
\caption{Why deepseek-r1 remains unusually strong on selected subsets.}
\label{tab:appendix_r1_capability}
\centering
\scriptsize
\setlength{\tabcolsep}{3.5pt}
\renewcommand{\arraystretch}{1.08}
\begin{tabular}{lp{18mm}p{22mm}p{31mm}p{31mm}}
\toprule
Diagnostic & R1 score & Stronger than & Rewards & Complementary loss \\
\midrule
P1 headline & 0.504 & gpt-5.4, deepseek-v3, qwen3.6-plus, claude-4.6-sonnet & balanced certification & still below qwen3-max and gemini-3.1-pro-preview \\
P2a first feasible & 29.8\% & gemini-3.1-pro-preview and all but qwen3.6-plus & initial anchoring & P2b headline still below gemini-3.1-pro-preview and claude-4.6-sonnet \\
P2c improvement rate & 0.7901 & effectively tied with gemini-3.1-pro-preview & sustained trajectory improvement & AUC still below gemini-3.1-pro-preview \\
P2b headline & 0.2413 & most of roster except gemini-3.1-pro-preview and claude-4.6-sonnet & closure under feedback & not the top endpoint model \\
P3-Success & 42.9\% & current point-estimate leader & corrupted-state recovery & RecoveryQuality still below gpt-5.4 \\
P4 Full Tau & 0.833 & mid-table only & relative structural ranking & clearly below gpt-5.4, qwen3.6-plus, and deepseek-v3 \\
\bottomrule
\end{tabular}
\end{table*}

\section{Stage-Gap Diagnostics}

The stage-optimal gap table quantifies how far each monolithic model remains from the oracle-routed stage-wise frontier. Gap is the mean of per-stage deficits after normalizing each stage by the corresponding oracle value. This table uses the frozen common-coverage split snapshot rather than the later expanded full-bank rows, so hunyuan-hy3-preview, mimo-v2.5-pro, and deepseek-v4-pro are excluded from this particular diagnostic. The point is system-level rather than leaderboard-level: if every single model remains far from the stage-wise envelope, then a verifier-gated, stage-routed engineering agent has headroom that a monolithic model choice cannot remove.

\begin{table*}[t]
\caption{Gap to the stage-wise envelope.}
\label{tab:appendix_stage_gap}
\centering
\scriptsize
\setlength{\tabcolsep}{5pt}
\renewcommand{\arraystretch}{1.08}
\begin{tabular}{p{38mm}p{34mm}p{44mm}}
\toprule
Model & Gap to stage-wise envelope & Largest deficit \\
\midrule
gemini-3.1-pro-preview & 0.062 & P3 (0.134) \\
deepseek-r1 & 0.141 & P2 (0.382) \\
qwen3-max & 0.214 & P2 (0.499) \\
gpt-5.4 & 0.232 & P2 (0.659) \\
deepseek-v3 & 0.291 & P2 (0.599) \\
\texttt{o4-mini} & 0.302 & P2 (0.603) \\
\texttt{qwen3.6-plus} & 0.327 & P1 (0.466) \\
claude-4.6-sonnet & 0.427 & P1 (0.640) \\
\texttt{llama-3.3-70b} & 0.546 & P3 (0.925) \\
\bottomrule
\end{tabular}
\end{table*}

\clearpage
\input{generated/appendix_extended_related_discussion.tex}

\clearpage
\section{CMA-ES Classical Optimization Baseline}

CMA-ES~\citep{hansen2023cma} is run on the full VEH P2 task bank under the same analytical oracle as the LLM agents, with population size 8 and initial sigma 0.15. The LLM budget (6 queries) and a relaxed budget (40 queries, 6.7$\times$ larger) are both tested for direct comparison with the model results in Table~\ref{tab:main_results} and the calibration baselines in Table~\ref{tab:appendix_calibration}.

\begin{table}[t]
\caption{CMA-ES P2 baseline results on the VEH 208-task bank. The 6-query budget matches the LLM agents; 40 queries tests the optimizer ceiling. The anchor-fixed heuristic (0 oracle queries) and Gemini (2.63 avg.\ queries) are shown for context.}
\label{tab:appendix_cmaes_p2}
\centering
\small
\begin{tabular}{lcccc}
\toprule
Method & Budget (queries) & Final feasible rate & Mean power ratio & Mean queries \\
\midrule
CMA-ES & 6 & 0.269 & 0.099 & 6.0 \\
CMA-ES & 40 & 0.578 & 0.305 & 40.0 \\
Anchor-fixed heuristic & 0 & 1.000 & 0.306 & 0.0 \\
model_B (LLM) & 6 & 0.962 & 0.390 & 2.63 \\
\bottomrule
\end{tabular}
\end{table}

CMA-ES under 6 queries underperforms most frontier LLM agents on P2. With 40 queries it matches the anchor-fixed heuristic (0.305 vs.\ 0.306) but remains below Gemini (0.390). This supports two complementary interpretations: (1) LLM agents carry physical priors from pretraining that accelerate feasibility-preserving search beyond what a pure numerical optimizer achieves under the same oracle budget; and (2) P2 is not solvable by a generic optimizer within the benchmark's query constraints, so the observed inter-model differences reflect genuine behavioral variation rather than a ceiling effect.


\clearpage
\section{Prompt Templates, Oracle Protocol, and Task Examples}
\label{sec:prompts}

\subsection{Evaluation Harness Protocol}

All model runs use a shared evaluation harness. Each task instance specifies the design variables, variable bounds, task objective, and query budget. The harness parses the first valid JSON object in each model response, normalizes field aliases (e.g., \texttt{R\_ohm} vs.\ \texttt{R\_value}), and applies a fixed retry policy: up to two retries for transient malformed or length-truncated outputs before counting the row as a parse failure. Iterative probes (P2, P3) run under bounded query budgets. P2 can stop on feasible closure; P3 preserves the full recovery trace. All runs use temperature 0.0, and provider-level default max output tokens (typically 4096--8192 depending on the provider). Run dates, provider endpoints, and model version strings are recorded in per-model run manifests released with the benchmark artifact.

\subsection{P1 Prompt Template (Triage)}

The P1 prompt presents the task specification and asks the model to output exactly one action.

\begin{quote}
\footnotesize
\textbf{System:} You are an engineering triage agent. Given a design specification, choose whether the task is feasible, infeasible, or underspecified. Do not output numeric design variables; output one action label as a single JSON object.

\textbf{User:}
\begin{verbatim}
{
  "task": "Evaluate whether the following VEH design specification
           is feasible, infeasible, or missing critical information.",
  "spec": {
    "beam_length_mm": [20, 80],
    "tip_mass_g": [0.5, 5.0],
    "substrate_thickness_um": [100, 500],
    "piezo_thickness_um": [50, 300],
    "target_frequency_hz": 50,
    "target_power_uw": 15,
    "excitation_acceleration_g": 0.5,
    "material": "PZT-5A",
    "substrate": "stainless_steel"
  },
  "allowed_actions": ["propose_design", "declare_infeasible",
                      "request_missing_info"]
}
\end{verbatim}

\textbf{Expected output format:}
\begin{verbatim}
{
  "action": "propose_design | declare_infeasible | request_missing_info",
  "reason": "<brief justification>",
  "missing_fields": ["spec.xxx", ...]
}
\end{verbatim}
\end{quote}

The \texttt{missing\_fields} array is only required for \texttt{request\_missing\_info}. P1 evaluator scoring is described in Appendix~\ref{sec:metric_formulas}.

\subsection{P2 Prompt Template (Iterative Repair)}

P2 runs as a multi-turn loop. Each turn presents the current design state and the oracle's constraint feedback.

\begin{quote}
\footnotesize
\textbf{System:} You are an engineering design agent. You are given a design task with coupled physical constraints. At each step, propose an improved design. The verifier will return feasibility and violation information. Your goal is to reach a feasible design within the query budget while maximizing the objective.

\textbf{User (turn 1):}
\begin{verbatim}
{
  "task": "Design a piezoelectric cantilever VEH meeting the
           following constraints.",
  "design_variables": ["beam_length_mm", "tip_mass_g",
       "substrate_thickness_um", "piezo_thickness_um",
       "beam_width_mm", "load_resistance_ohm"],
  "variable_bounds": {
    "beam_length_mm": [20, 80],
    "tip_mass_g": [0.5, 5.0],
    "substrate_thickness_um": [100, 500],
    "piezo_thickness_um": [50, 300],
    "beam_width_mm": [5, 30],
    "load_resistance_ohm": [1000, 100000]
  },
  "initial_design": {
    "beam_length_mm": 50, "tip_mass_g": 2.0,
    "substrate_thickness_um": 300, "piezo_thickness_um": 150,
    "beam_width_mm": 15, "load_resistance_ohm": 10000
  },
  "objective": "maximize output power (uw)",
  "query_budget": 6
}
\end{verbatim}

\textbf{Oracle response (example, after turn 1):}
\begin{verbatim}
{
  "feasible": false,
  "metrics": {
    "resonant_freq_hz": 72.3,
    "load_power_uw": 3.2,
    "tip_stress_mpa": 12.4,
    "tip_disp_mm": 0.8
  },
  "violations": [
    {"constraint": "frequency_target", "value": 72.3,
     "target": 50.0, "tolerance": 0.05,
     "message": "frequency too high; increase mass or length"}
  ],
  "total_violation": 8.47,
  "oracle_feedback": "frequency=72.3 Hz exceeds [47.5, 52.5] Hz
      target; tip stress=12.4 MPa within limit."
}
\end{verbatim}

\textbf{Model output format (each turn):}
\begin{verbatim}
{
  "design": {
    "beam_length_mm": 55.0,
    "tip_mass_g": 2.5,
    "substrate_thickness_um": 280,
    "piezo_thickness_um": 150,
    "beam_width_mm": 15,
    "load_resistance_ohm": 10000
  },
  "action": "continue",
  "reason": "Increased beam length and tip mass to reduce
             resonant frequency toward 50 Hz target."
}
\end{verbatim}
\end{quote}

\subsection{P3 Prompt Template (Post-Trap Recovery)}

P3 is structurally identical to P2 but pre-loads a corrupted trajectory into the conversation context. The corruption varies by trap type: unit-flip, wrong-formula-direction, false-feasibility, topology-trap, verifier-ignored, or progressive-contamination. The model receives the full pseudo-history as if it were prior conversation turns.

For the state-summary intervention, the raw trajectory is replaced with a deterministic verifier-authored summary containing only already-observed trusted fields: current step, latest proposal, latest verifier state, best-so-far feasible proposal (if any), short-horizon objective and violation trends, and whether new violations were introduced. The oracle evaluator is unchanged.

\subsection{P4 Prompt Template (Policy-Conditioned Ranking)}

\begin{quote}
\footnotesize
\textbf{System:} You are an engineering evaluator. Given a set of oracle-feasible design candidates and a stated ranking policy, rank the candidates according to the policy. Do not modify the designs or propose new ones.

\textbf{User:}
\begin{verbatim}
{
  "task": "Rank the following feasible VEH designs according to
           the stated policy.",
  "policy": "Prioritize output power, then tip-stress margin,
             then component cost. All candidates are feasible.",
  "candidates": [
    {"id": "A", "design": {...}},
    {"id": "B", "design": {...}},
    {"id": "C", "design": {...}},
    {"id": "D", "design": {...}},
    {"id": "E", "design": {...}}
  ]
}
\end{verbatim}

\textbf{Expected output:}
\begin{verbatim}
{
  "ranking": ["B", "A", "D", "E", "C"],
  "reason": "B has highest power; A second on power with better
             stress margin than D; E and C follow by power."
}
\end{verbatim}
\end{quote}

\subsection{Oracle Specification (VEH Domain)}

The VEH oracle follows a single-mode analytical model for a piezoelectric cantilever beam with tip mass. For a design vector $\mathbf{x} = (L, m_t, h_s, h_p, w, R_L)$ and excitation acceleration $a$ at frequency $f_{\text{exc}}$, the oracle computes:

\begin{align}
f_r &= \frac{1}{2\pi}\sqrt{\frac{k_{\text{eff}}}{m_{\text{eff}}}} \quad\text{(resonant frequency)} \\
P_{\text{out}} &= \frac{V_{\text{out}}^2}{R_L} \quad\text{(output power)} \\
\sigma_{\max} &= \frac{M_{\max} \cdot y_{\max}}{I} \quad\text{(tip stress)} \\
\delta_{\max} &= \frac{F_{\text{eff}} L^3}{3 E I} \quad\text{(tip displacement)}
\end{align}

where $k_{\text{eff}}$, $m_{\text{eff}}$, $V_{\text{out}}$, $M_{\max}$, $y_{\max}$, \(I\), and \(F_{\text{eff}}\) are derived from $\mathbf{x}$ using the released oracle implementation. The oracle returns feasible status, metrics, violations, and verifier feedback. A design is feasible iff $|f_r-f_{\text{exc}}|/f_{\text{exc}}\le\texttt{freq\_tol}$, $\sigma_{\max}\le\texttt{stress\_limit}$, and $\delta_{\max}\le\texttt{disp\_limit}$. The release includes the oracle code, task manifests, and validation artifacts needed to audit these computations.

\subsection{Example Task Walkthrough (VEH P2)}

We illustrate a complete P2 trajectory for a single task to make the evaluation protocol concrete.

\textbf{Task.} Design a PZT-5A cantilever VEH with target frequency 50 Hz, target power $\ge 15\,\mu$W, excitation 0.5 g, stress limit 50 MPa, displacement limit 2 mm. Budget: 6 queries. Initial design: $L = 50$ mm, $m_t = 2.0$ g, $h_s = 300\,\mu$m, $h_p = 150\,\mu$m, $w = 15$ mm, $R_L = 10$ k$\Omega$.

\textbf{Step 1.} Model proposes $\{L: 55, m_t: 2.5, h_s: 280, h_p: 150, w: 15, R_L: 10000\}$. Oracle returns: $f_r = 63.8$ Hz (violation: too high), $P_{\text{out}} = 4.1\,\mu$W, stress 11.2 MPa, displacement 0.7 mm. Feasible: false. Total violation: 5.24.

\textbf{Step 2.} Model proposes $\{L: 65, m_t: 3.0, h_s: 250, h_p: 120, w: 15, R_L: 10000\}$. Oracle returns: $f_r = 48.7$ Hz (within $[47.5, 52.5]$), $P_{\text{out}} = 9.8\,\mu$W, stress 9.1 MPa, displacement 0.8 mm. Feasible: true. Violation: 0. Objective: $9.8/15 = 0.653$ (relative to BKF). End of trajectory (feasible closure).

This task is considered a P2 success. The model used 2 queries, achieved feasible closure, and the final feasible power ratio is 0.653.

\subsection{Example Task Walkthrough (Circuit P3 Dual-Trap)}

\textbf{Task.} RC low-pass filter, target $f_c = 1$ kHz $\pm 2\%$, source current limit 0.5 mA, $R \in [1\text{k}\Omega, 100\text{k}\Omega]$, $C \in [1\text{nF}, 1\mu\text{F}]$, $V_{\text{in}} = 5$ V. Corrupted history (phase-1 trap): prior steps repeatedly increased $C$ under the false assumption that $f_c$ was too low; actual $f_c = 3.18$ kHz (too high). Current bad design: $R = 5\text{k}\Omega$, $C = 100$ nF. Budget: 5 queries.

\textbf{Step 1.} Model recognizes the contradiction between oracle feedback (``$f_c$ too high'') and corrupted history (``$f_c$ too low''). Proposes $\{R: 20000, C: 7.96\text{nF}\}$. Oracle: $f_c = 1000$ Hz (feasible), $I_{\text{source}} = 0.25$ mA (feasible). Phase 1 escaped.

\textbf{Step 2 (phase-2 trap).} With $R = 20\text{k}\Omega$, the model's edit has triggered the near-boundary margin: a second constraint on component tolerance (previously dormant) now requires $C \ge 1.5$ nF. Oracle returns a new violation. Model must adjust $R$ downward slightly and $C$ upward.

\textbf{Step 3.} Model proposes $\{R: 15000, C: 10.6\text{nF}\}$. Oracle: $f_c = 1001$ Hz (feasible), $I_{\text{source}} = 0.33$ mA (feasible). All constraints satisfied. Recovery success: true. Cascade: false.

This task demonstrates the progressive-dual-trap mechanism: the model must escape the primary corrupted state (phase 1) and then handle a secondary coupled constraint triggered by the escape move itself (phase 2).

\subsection{Run Configuration}

All model runs share the following configuration unless noted otherwise in per-model manifests:

\begin{itemize}[leftmargin=5mm,itemsep=1mm,topsep=1mm]
    \item Temperature: 0.0 (deterministic decoding).
    \item Max output tokens: provider default (4096--8192).
    \item Retry policy: up to 2 retries on parse failure; row counted as parse error after 3 consecutive failures.
    \item P2 stop rule: trajectory terminates on feasible closure or budget exhaustion.
    \item P3 stop rule: trajectory runs full budget; no early stop on feasibility.
    \item Provider endpoints and model version strings are recorded in per-model run manifests.
    \item Core roster run dates: April 2026. Thinking-model extension run dates: April 25--27, 2026. Circuit audit run dates: April 28--30, 2026.
\end{itemize}


\clearpage
\section{Representative Failure Cases}
\label{sec:failure_cases}

The following notes give one auditable failure case per probe. They are not additional metrics; they make the response-control interpretation concrete and show what the probe-specific scores count as failure. We use prose rather than a table because the cases contain long task identifiers and multi-sentence readouts.

\paragraph{P1: optimistic over-proposal.}
In audit case \path{p1::dev::infeasible_hard_conflict::0000::s42}, the gold action is \texttt{declare\_infeasible}, but the model chooses \texttt{propose\_design}. The verifier-side label makes the failure observable: continuation is counted as a spurious propose and unsafe entry into search. This is an action-prior failure. P1 is therefore not rewarding raw willingness to generate a design; it rewards disciplined entry control.

\paragraph{P2: over-edit / wandering.}
In the macro-unimorph tip-mass repair case derived from DOI 10.3390/mi14020421 (case 0016), the repair trace makes large moves across violation families instead of preserving the feasible neighborhood. The oracle trace shows oscillation and failure to re-enter feasibility within budget; over-edit and low directed-repair scores capture this behavior. This is an edit-style failure: the model can react to feedback, but the reaction is too global and can chase one constraint while breaking another.

\paragraph{P3: continuity trap.}
In audit case \path{p3::dev::0093::0000}, the exposed history has already moved substrate thickness in a harmful direction. Continuity-biased runs keep treating that corrupted history as trusted state. P3 separates this into escape, cascade, dead-budget, and final-success fields: merely noticing the trap is not enough if the model then cascades or fails to return to a feasible design. This is why the state-summary intervention is diagnostic; it tests whether verifier-authored state isolation reduces raw-history continuity bias.

\paragraph{P4: policy mismatch.}
In audit case \path{diagbench::test_id::p4::0032}, all candidates are physically feasible, but the model ranks a local trade-off neighborhood as \(B>C>E>A>D\) while the oracle policy ranks \(C>E>D>A>B\). The resulting full Kendall tau is \(-1.0\) and the dominance-violation rate is \(1.0\). The error is therefore not search failure. It is preference-execution failure: the model does not follow the policy-conditioned ordering over feasible designs.

Together, these cases clarify why DiagBench reports stage-specific metrics instead of a single endpoint score. A P1 over-proposal can look productive but is unsafe triage; a P2 over-edit can contain valid engineering arithmetic but still destroy feasibility; a P3 continuity trap can escape one violation and cascade into another; and a P4 policy mismatch can occur after the physical verifier has already accepted every candidate.


\clearpage
\section{Selection--Generation Isomorphism Audit}
\label{sec:appendix_isomorphic}

The following tables document the frozen cleanup protocol behind the SG-gap analysis in the main text. Residual failures after one cleanup pass and at most one targeted rerun are reported as final outcomes rather than recursively cleaned away.

\subsection{Selection--Generation Isomorphism Audit}

This audit bank tests whether the same underlying feasible design remains accessible when the response regime changes. We construct 46 probe groups from the hard/medium P2--P4 overlap and instantiate each group in three matched forms: \textbf{A-selection} (pick the one feasible candidate from a five-way pool), \textbf{B-generation} (synthesize a feasible candidate from a near-feasible seed), and \textbf{C-completion} (fill in two masked variables of a gold feasible design). We freeze the protocol at one cleanup pass plus at most one targeted rerun for malformed outputs; residual failures after that pass are reported as final failures rather than recursively cleaned away.

\begin{table*}[t]
\caption{Isomorphic probe results after one cleanup pass and at most one targeted rerun. Success rates are group-level percentages over 46 probe groups. SG-gap is A-selection minus B-generation. Gemini remains partially parse-limited only on B-generation; all other model/form cells are parse-clean after the frozen cleanup protocol.}
\label{tab:appendix_isomorphic_probe}
\centering
\scriptsize
\setlength{\tabcolsep}{3pt}
\renewcommand{\arraystretch}{1.08}
\begin{tabular}{p{30mm}ccccccc}
\toprule
Model & A-sel. & B-gen. & C-comp. & SG-gap & Parse-B & Lat. (s) & Bad rows \\
\midrule
model_E & 41.3\% & 4.3\% & 17.4\% & 37.0 pts & 100.0\% & 18.1 & 0 \\
model_C & 32.6\% & 6.5\% & 21.7\% & 26.1 pts & 100.0\% & 2.4 & 0 \\
model_B & 65.2\% & 4.3\% & 30.4\% & 60.9 pts & 41.3\% & 29.8 & 27 \\
\texttt{model_G} & 52.2\% & 8.7\% & 15.2\% & 43.5 pts & 100.0\% & 70.0 & 0 \\
\bottomrule
\end{tabular}
\end{table*}

\begin{table*}[t]
\caption{Protocol log for the isomorphic probe reruns. Residual bad rows after the first targeted rerun are treated as part of the measured behavior, not recursively retried until clean.}
\label{tab:appendix_isomorphic_protocol}
\centering
\scriptsize
\setlength{\tabcolsep}{3pt}
\renewcommand{\arraystretch}{1.08}
\begin{tabularx}{\textwidth}{p{28mm}Y p{28mm}Y}
\toprule
Model & First-pass issue & Final residual after one rerun & Readout \\
\midrule
model_C & no malformed-output cleanup needed & 0 bad rows & stable baseline; all A/B/C cells remain parse-clean \\
model_E & transient transport failures in the first pass, cleaned once & 0 bad rows & clean selector after one repair pass; low B-generation remains a capability result, not a format artifact \\
model_B & first pass mixed timeouts and malformed JSON & 27 bad rows, all B-generation & residual failures are no longer timeout-dominated; they remain half-written JSON under the hardest six-variable generation form \\
\texttt{model_G} & four timeout rows in the first pass & 0 bad rows & fully clean after one repair pass, but still by far the slowest model in the audit \\
\bottomrule
\end{tabularx}
\end{table*}


\clearpage
\clearpage
\section{Tier~3 Controlled-Prompt Audit --- Full Results}
\label{sec:tier3_full}

The main text (Section~\ref{sec:experiments}) reports the four primary confirmatory cells with targeted-minus-neutral paired bootstrap deltas. This appendix provides the complete per-condition results, the C4 per-model split, and per-task pattern decomposition that support the interpretation advanced in the main text.

\subsection{Design and Execution}

The experiment was fully pre-registered. Four confirmatory cells pair known deficit models with their weak stages: C1 (model_C, P2), C2 (model_E, P3), C3 (model_B, P4), and C4 (model_C and model_E, P1). Each cell samples 36--60 tasks under frozen seed 42 with subtype-level stratification. Five prompt conditions are tested per cell: \texttt{default} (the standard benchmark prompt), \texttt{neutral} (generic structured execution without boundary-specific diagnosis), \texttt{targeted} (boundary-specific diagnostic prompt designed to repair the known deficit), \texttt{wrong-boundary} (a mismatched diagnostic prompt from a different stage), and for C2, a \texttt{state-summary-plus} variant stacking targeted with default. All prompts were frozen before execution; no iterative refinement occurred. The analysis uses 10,000 paired task-level bootstrap resamples with 95\% confidence intervals.

\subsection{Full Per-Condition Results}

\begin{table}[t]
\caption{Tier~3 Phase~A complete per-condition results. All metrics are primary confirmatory endpoints per cell.}
\label{tab:tier3_all_conditions}
\centering
\scriptsize
\setlength{\tabcolsep}{3.5pt}
\renewcommand{\arraystretch}{1.08}
\begin{tabular}{lllcccc}
\toprule
Cell & Model & Condition & $n$ & Primary metric & Secondary & Parse \\
\midrule
C1 & model_C & default   & 40 & 0.150 feasible & 61.0 mean obj & 0.000 \\
   &         & neutral   & 40 & 0.175 feasible & 4.3 mean obj  & 0.000 \\
   &         & targeted  & 40 & 0.150 feasible & 461.8 mean obj & 0.000 \\
   &         & wrong     & 40 & 0.075 feasible & 2.1 mean obj  & 0.000 \\
\midrule
C2 & model_E & default   & 36 & 0.167 success & 0.694 cascade & 0.000 \\
   &            & neutral   & 36 & 0.389 success & 0.528 cascade & 0.000 \\
   &            & targeted  & 36 & 0.194 success & 0.778 cascade & 0.000 \\
   &            & wrong     & 36 & 0.222 success & 0.694 cascade & 0.000 \\
   &            & s+        & 36 & 0.222 success & 0.750 cascade & 0.000 \\
\midrule
C3 & Gemini-3.1 & default   & 50 & 0.692 $\tau$ & 0.320 exact & 0.000 \\
   &            & neutral   & 50 & 0.692 $\tau$ & 0.320 exact & 0.000 \\
   &            & targeted  & 50 & 0.708 $\tau$ & 0.320 exact & 0.000 \\
   &            & wrong     & 50 & 0.700 $\tau$ & 0.300 exact & 0.000 \\
\midrule
C4 & model_E & default   & 60 & 0.650 acc & --- & 0.050 \\
   &            & neutral   & 60 & 0.617 acc & --- & 0.050 \\
   &            & targeted  & 60 & 0.567 acc & --- & 0.183 \\
   &            & wrong     & 60 & 0.633 acc & --- & 0.050 \\
\cmidrule{2-7}
   & model_C    & default   & 60 & 0.433 acc & --- & 0.000 \\
   &            & neutral   & 60 & 0.400 acc & --- & 0.000 \\
   &            & targeted  & 60 & 0.467 acc & --- & 0.000 \\
   &            & wrong     & 60 & 0.400 acc & --- & 0.000 \\
\bottomrule
\end{tabular}
\end{table}

Table~\ref{tab:tier3_all_conditions} reports the complete per-condition primary metric. Three patterns merit attention beyond the main-text targeted-minus-neutral deltas:

\begin{enumerate}[leftmargin=5mm,itemsep=1mm,topsep=1mm]
    \item \textbf{C2 neutral improvement.} The generic structured neutral prompt lifts Claude P3 success from 0.167 to 0.389 (95\% CI for neutral$-$default: [0.083, 0.361]), while reducing cascade from 0.694 to 0.528. This is the largest single-condition effect in the experiment and suggests that Claude benefits from structured execution guidance but not from skeptical-reset framing.
    \item \textbf{C2 cascade--success inversion.} Cascade rate under targeted (0.778) is the highest across all C2 conditions, and targeted is the only condition where cascade exceeds the default rate. Per-task analysis reveals four tasks with cascade exclusively under targeted, consistent with the interpretation that the skeptical-reset prompt induces overcorrection.
    \item \textbf{C4 parse-error interaction.} Claude's parse-error rate under targeted P1 (0.183) is more than triple the rate under other conditions (0.050). Parse errors concentrate in \texttt{infeasible\_margin} tasks (4/7) and \texttt{declare\_infeasible} gold labels (5/14). Excluding parse errors, Claude's classification accuracy under targeted is 0.694, exceeding neutral (0.617), indicating that the targeted prompt content aids reasoning while its length impairs output-format compliance.
\end{enumerate}

\subsection{Per-Task Decomposition}

\paragraph{C2 per-task transitions.} Of 36 tasks, 20 fail under all three main conditions; 7 succeed under neutral but fail under default and targeted; 4 succeed under all conditions. The neutral$\to$targeted transition loses 9 successes and gains 2. Cascade under targeted increases in 10 tasks relative to neutral and decreases in only 1.

\paragraph{C3 per-task decomposition.} The +0.016 mean $\tau$ gain is driven entirely by one task whose $\tau$ moves from $-$0.4 (negatively correlated ranking) under neutral to +0.4 under targeted; the remaining 49 tasks show identical $\tau$ values across both conditions. Top-1 accuracy improves on 2 tasks (94\%$\to$100\%). The targeted prompt thus eliminates a single ranking failure rather than systematically improving order quality.

\paragraph{C1 subtype pattern.} model_C never solves \texttt{boundary\_binding} tasks under any condition (0/40). The \texttt{paper\_like} subtype accounts for most feasible cases across all conditions (19/22 total feasible designs). This subtype specificity, together with the uniformly low feasible rates (0.075--0.175), indicates a capability floor that prompt variation does not shift.

\subsection{Interpretive Summary}

The four-cell audit supports the following interpretation, which we advance in the main text: behavioral fingerprints identified by Tier~1 profiles are not shallow prompt artifacts. If they were, targeted boundary-specific prompts should repair them. Instead, C1 shows no effect, C2 shows significant reversal, C3 shows a one-task marginal effect, and C4 shows model-dependent interaction. The resilience of stage dissociation to prompt-level repair reinforces the benchmark's construct validity: what DiagBench measures is durable prior-boundary compatibility, not prompt sensitivity.


\clearpage
\section{Circuit Audit: Cross-Domain Construct Validity}

This appendix reports the final circuit audit described in Section~\ref{sec:experiments}. The audit is a construct-validity check, not an independent circuit benchmark: it asks whether the same P1--P4 diagnostic regimes remain discriminative under a second closed-form engineering domain. The released task bank hardens P1 with near-boundary infeasible and missing-information cases, hardens P2 with dual-constraint edits where fixing one violation can break another, and uses matched P3/P4 recovery and ranking banks to test corrupted-state recovery and policy-conditioned selection.

\Needspace{12\baselineskip}
\begin{center}
\captionof{table}{Circuit audit task inventory. The task bank changes the physics while preserving the P1--P4 decision boundaries.}
\label{tab:appendix_circuit_inventory}
\scriptsize
\setlength{\tabcolsep}{3pt}
\renewcommand{\arraystretch}{1.08}
\begin{tabularx}{\textwidth}{lcYYY}
\toprule
Probe & n & Source & Hardening & Validation target \\
\midrule
P1 & 32 & final circuit audit & Near-boundary infeasible cases; lower raw \texttt{propose\_design} prior; subtype-balanced entry decisions & Action-prior transfer \\
P2 & 32 & final circuit audit & Dual constraints; repair directions can conflict; objective-preserving feasible repair is rewarded & Edit style and feedback obedience \\
P3 & 18 & final circuit audit & Progressive dual traps; escape can trigger a second violation & State-trust policy and trap sensitivity \\
P4 & 24 & final circuit audit & Policy-flip ranking tasks under feasible candidate pools & Preference execution \\
\bottomrule
\end{tabularx}
\end{center}

\paragraph{Why P1-Composite is the circuit P1 headline.}
The circuit P1 bank follows the VEH P1 hardening principle: raw accuracy is too easy to inflate when the action distribution contains many feasible proposals. The bank therefore lowers the trivial \texttt{propose\_design} prior and adds near-boundary infeasible/request cases. P1-Composite is used as the headline because it jointly rewards correct entry actions and penalizes spurious proposals, unsafe proposals, missed missing-information cases, and missed infeasibility cases. This makes P1 an entry-discipline metric rather than a willingness-to-generate metric.

\Needspace{12\baselineskip}
\begin{center}
\captionof{table}{Circuit P1 action triage. P1-Comp. is the headline; abbreviations: Spur.=spurious propose, Unsafe=unsafe propose, Req.=request recall, Inf.=infeasible recall.}
\label{tab:appendix_circuit_p1}
\scriptsize
\setlength{\tabcolsep}{3.5pt}
\renewcommand{\arraystretch}{1.08}
\begin{tabular}{lccccccccc}
\toprule
Model & n & P1-Comp. & Acc. & F1 & Spur. & Unsafe & Req. & Inf. & Parse \\
\midrule
qwen3-max & 32 & 0.971 & 0.969 & 0.965 & 0.000 & 0.000 & 1.000 & 1.000 & 0.000 \\
hunyuan-hy3-preview & 32 & 0.944 & 0.938 & 0.933 & 0.000 & 0.000 & 1.000 & 1.000 & 0.000 \\
deepseek-r1 & 32 & \textbf{0.974} & 0.969 & 0.971 & 0.000 & 0.000 & 1.000 & 1.000 & 0.000 \\
gemini-3.1-pro-preview & 32 & \textbf{0.974} & 0.969 & 0.971 & 0.000 & 0.000 & 1.000 & 1.000 & 0.000 \\
gpt-5.4 & 32 & 0.923 & 0.906 & 0.917 & 0.000 & 0.000 & 1.000 & 1.000 & 0.000 \\
claude-4.6-sonnet & 32 & 0.932 & 0.938 & 0.940 & 0.031 & 0.031 & 1.000 & 0.875 & 0.000 \\
\bottomrule
\end{tabular}
\end{center}

\Needspace{12\baselineskip}
\begin{center}
\captionof{table}{Circuit P2 iterative repair. P2b is the final feasible objective score; abbreviations: Feas.=final feasible rate, Dir.=directed repair, O-edit=over-edit.}
\label{tab:appendix_circuit_p2}
\scriptsize
\setlength{\tabcolsep}{4pt}
\renewcommand{\arraystretch}{1.08}
\begin{tabular}{lcccccc}
\toprule
Model & n & Feas. & P2b & Dir. & O-edit & Parse \\
\midrule
qwen3-max & 32 & 0.906 & 0.494 & 0.938 & 0.400 & 0.000 \\
hunyuan-hy3-preview & 32 & 0.938 & 0.512 & 0.914 & 0.771 & 0.000 \\
deepseek-r1 & 32 & 0.969 & 0.594 & \textbf{1.000} & 0.605 & 0.000 \\
gemini-3.1-pro-preview & 32 & \textbf{1.000} & \textbf{0.662} & 0.969 & 0.938 & 0.000 \\
gpt-5.4 & 32 & 0.875 & 0.457 & 0.883 & 0.489 & 0.000 \\
claude-4.6-sonnet & 32 & 0.906 & 0.487 & 0.914 & 0.646 & 0.000 \\
\bottomrule
\end{tabular}
\end{center}

\Needspace{12\baselineskip}
\begin{center}
\captionof{table}{Circuit P3 corrupted-state recovery retained from the original bank. Abbreviations: Esc.=escape, Replan=explicit replan, Reset=history reset, Casc.=cascade, Dead=dead budget, Succ.=final success.}
\label{tab:appendix_circuit_p3}
\scriptsize
\setlength{\tabcolsep}{3.5pt}
\renewcommand{\arraystretch}{1.08}
\begin{tabular}{lccccccccc}
\toprule
Model & n & Esc. & Replan & Reset & Casc. & Dead & Succ. & Rec. q. & Parse \\
\midrule
qwen3-max & 18 & 0.944 & 0.000 & 0.000 & 0.471 & 0.000 & 0.778 & 0.366 & 0.000 \\
gemini-3.1-pro-preview & 18 & 1.000 & 0.000 & 0.000 & \textbf{0.000} & 0.000 & \textbf{1.000} & \textbf{0.611} & 0.000 \\
gpt-5.4 & 18 & 0.944 & 0.111 & 0.111 & 0.176 & 0.056 & 0.778 & 0.401 & 0.000 \\
deepseek-r1 & 18 & 0.778 & 0.000 & 0.000 & \textbf{0.000} & 0.000 & 0.722 & 0.314 & 0.000 \\
hunyuan-hy3-preview & 18 & 0.889 & 0.000 & 0.000 & 0.062 & 0.000 & 0.778 & 0.364 & 0.000 \\
claude-4.6-sonnet & 18 & 1.000 & 0.000 & 0.000 & \textbf{0.000} & 0.000 & \textbf{1.000} & 0.568 & 0.000 \\
\bottomrule
\end{tabular}
\end{center}

\Needspace{12\baselineskip}
\begin{center}
\captionof{table}{Circuit P4 policy-conditioned ranking. Abbreviations: Top2=set agreement for the top two candidates; Flip=policy-flip accuracy.}
\label{tab:appendix_circuit_p4}
\scriptsize
\setlength{\tabcolsep}{3.5pt}
\renewcommand{\arraystretch}{1.08}
\begin{tabular}{lccccccccc}
\toprule
Model & n & Tau & Exact & Top1 & Top2 & Pair & Flip & BARS & Parse \\
\midrule
qwen3-max & 24 & 0.608 & 0.375 & 0.583 & 0.583 & 0.804 & 0.769 & 0.710 & 0.000 \\
gemini-3.1-pro-preview & 24 & 0.475 & 0.250 & 0.500 & 0.500 & 0.738 & 0.679 & 0.625 & 0.000 \\
gpt-5.4 & 24 & \textbf{0.625} & 0.375 & 0.583 & 0.625 & 0.812 & \textbf{0.801} & \textbf{0.722} & 0.000 \\
deepseek-r1 & 24 & 0.542 & 0.292 & 0.583 & 0.500 & 0.771 & 0.716 & 0.661 & 0.000 \\
hunyuan-hy3-preview & 24 & 0.492 & 0.208 & 0.458 & 0.417 & 0.746 & 0.647 & 0.614 & 0.000 \\
claude-4.6-sonnet & 24 & 0.608 & 0.375 & 0.542 & 0.500 & 0.804 & 0.759 & 0.707 & 0.000 \\
\bottomrule
\end{tabular}
\end{center}

\Needspace{10\baselineskip}
\begin{center}
\captionof{table}{Circuit response-control profile scores. Action/edit scores use the updated P1/P2 audit; state/preference scores are retained from the original P3/P4 banks.}
\label{tab:appendix_circuit_profile}
\scriptsize
\setlength{\tabcolsep}{5pt}
\renewcommand{\arraystretch}{1.08}
\begin{tabular}{lcccc}
\toprule
Model & Action & Edit & State & Preference \\
\midrule
qwen3-max & 0.983 & 0.877 & 0.650 & 0.706 \\
hunyuan-hy3-preview & 0.968 & 0.792 & 0.721 & 0.612 \\
deepseek-r1 & \textbf{0.985} & 0.579 & 0.700 & 0.672 \\
gemini-3.1-pro-preview & \textbf{0.985} & 0.766 & \textbf{0.800} & 0.633 \\
gpt-5.4 & 0.956 & 0.830 & 0.720 & \textbf{0.714} \\
claude-4.6-sonnet & 0.962 & 0.815 & \textbf{0.800} & 0.696 \\
\bottomrule
\end{tabular}
\end{center}

Table~\ref{tab:appendix_cross_domain_patterns} summarizes the signature patterns that reproduce across VEH and circuit domains.

\Needspace{14\baselineskip}
\begin{center}
\captionof{table}{Cross-domain pattern reproduction. The circuit audit checks whether the same diagnostic reversals appear under different physics.}
\label{tab:appendix_cross_domain_patterns}
\scriptsize
\setlength{\tabcolsep}{3pt}
\renewcommand{\arraystretch}{1.08}
\begin{tabularx}{\textwidth}{p{31mm}YY}
\toprule
Pattern & VEH observation & Circuit observation \\
\midrule
gpt-5.4 P4 lead / weaker repair & P4 Tau 0.887 (1st); P2 ratio 0.133 & retained P4 Tau 0.625 (1st); circuit P2b 0.457 (6/6) \\
gemini-3.1-pro-preview P2 lead / weak P4 & P2 ratio 0.390 (1st); P4 Tau 0.824 (7/12) & circuit P2b 0.662 (1st); retained P4 Tau 0.475 (6/6) \\
No single model dominates & qwen3-max P1; gemini-3.1-pro-preview P2; hunyuan-hy3-preview P3; gpt-5.4 P4 & gemini-3.1-pro-preview/deepseek-r1 P1; gemini-3.1-pro-preview P2; gemini-3.1-pro-preview/claude-4.6-sonnet retained P3; gpt-5.4 retained P4 \\
P3 trap-geometry sensitivity & Multi-step contamination penalizes continuity-biased models (gemini-3.1-pro-preview cascade 0.507 in VEH) & Formula-level circuit traps give gemini-3.1-pro-preview zero cascade, while qwen3-max cascades at 0.471; this is a trap-geometry contrast, not a claim that the same model-level cascade mechanism transfers unchanged \\
Profile diagonal direction & Action prior tracks P1; edit style tracks P2; preference execution tracks P4 & P1-Composite, P2b, and P3 success remain separable under the circuit oracle \\
\bottomrule
\end{tabularx}
\end{center}

\paragraph{What the circuit audit proves and does not prove.}
The audit supports a narrow construct-validity claim: when the physics changes from VEH to closed-form circuit design, the P1--P4 scaffold still produces separable action, repair, recovery, and ranking behaviors. It also shows that response-control reversals such as gemini-3.1-pro-preview's repair/ranking split and gpt-5.4's ranking/repair split are not artifacts of VEH-specific formulas. It does \emph{not} establish stable circuit-domain model rankings, cover modern circuit design broadly, or replace a full circuit benchmark; the task count and family coverage are intentionally pilot-scale.

